% Algebra 2 Lesson 2-7 Lesson Quiz, short form (fast-lesson draft 2026-10-31, re-cut toward the Topic 2 test;
% docs/a2-fast-lesson-spec.md section 6). Built from the "2-7 Lesson Quiz: Linear-Quadratic Systems" block of a2_2-7_TE.tex.
% The form keeps the Savvas item numbers (spec section 6): items 1 (Example 1), 2 (bonus) and 5 (Example 5).
% Savvas items 3 and 4 (Examples 3 and 4, cut in the fast packet) are left off and listed in the key.
\documentclass[11pt]{article}
\usepackage[margin=0.7in]{geometry}
\usepackage{amsmath,amssymb}
\usepackage{enumitem}
\usepackage{tikz}
\usepackage{fancyhdr}
\usepackage{parskip}
\setlist[enumerate]{leftmargin=*, itemsep=6pt}
\pagestyle{fancy}
\fancyhf{}
\fancyfoot[L]{\small\textit{en}Vision\textsuperscript{\textregistered} Algebra 2}
\fancyfoot[C]{\small 2-7 Lesson Quiz, short form}
\fancyfoot[R]{\small\thepage}
\renewcommand{\headrulewidth}{0pt}
\newcommand{\choice}[1]{\tikz[baseline=-0.6ex]\node[draw, circle, inner sep=0.6pt, minimum size=3.3mm, font=\scriptsize\bfseries] {#1};\;}

\begin{document}

\noindent\textbf{Name}\; \rule{6cm}{0.4pt} \hfill Period \rule{1.2cm}{0.4pt} \quad Date \rule{2.2cm}{0.4pt}

\vspace{4pt}\noindent\rule{\textwidth}{0.6pt}

\vspace{2pt}\noindent{\LARGE\textbf{2-7 \; Lesson Quiz: Linear-Quadratic Systems}} \hfill {\small short form}

\vspace{4pt}\noindent\rule{\textwidth}{0.6pt}

\vspace{2pt}\noindent{\small This form has Savvas items 1, 2 and 5. Items 3 and 4 are on the full Savvas quiz only. Item 5 needs a graphing calculator (or Desmos).}

\vspace{4pt}
\begin{minipage}[t]{0.46\textwidth}
\textbf{1.} Determine the number of real solutions of the system $y=x^2-3$, $y=-2x+4$.

\vspace{4pt}
\begin{itemize}[leftmargin=1.2em, itemsep=5pt]
\item[] \choice{A} 2
\item[] \choice{B} 1
\item[] \choice{C} 3
\item[] \choice{D} 0
\end{itemize}
\end{minipage}\hfill
\begin{minipage}[t]{0.50\textwidth}
Show how you know (a sketch or your algebra):

\vspace{16pt}\rule{\linewidth}{0.4pt}

\vspace{16pt}\rule{\linewidth}{0.4pt}

\vspace{16pt}\rule{\linewidth}{0.4pt}

\vspace{16pt}\rule{\linewidth}{0.4pt}
\end{minipage}

\vspace{12pt}\noindent\rule{\textwidth}{0.4pt}

\vspace{4pt}
\begin{minipage}[t]{0.46\textwidth}
\textbf{2.} Use substitution to solve the system $y=-\frac{1}{2}x^2$, $y=x-4$.

\vspace{4pt}
\begin{itemize}[leftmargin=1.2em, itemsep=5pt]
\item[] \choice{A} $-4$ and $2$
\item[] \choice{B} $(4,-8)$ and $(-2,-2)$
\item[] \choice{C} $(-4,-8)$ and $(2,-2)$
\item[] \choice{D} $(0,0)$ and $(0,-4)$
\end{itemize}
\end{minipage}\hfill
\begin{minipage}[t]{0.50\textwidth}
Show your substitution:

\vspace{16pt}\rule{\linewidth}{0.4pt}

\vspace{16pt}\rule{\linewidth}{0.4pt}

\vspace{16pt}\rule{\linewidth}{0.4pt}

\vspace{16pt}\rule{\linewidth}{0.4pt}
\end{minipage}

\vspace{12pt}\noindent\rule{\textwidth}{0.4pt}

\vspace{4pt}
\textbf{5.} Solve the equation $2x^2+5x-8=\frac{5}{2}x+20$ by writing a linear-quadratic system and solving using the intersection feature of a graphing calculator. Round to the nearest tenth.

\vspace{8pt}
The system I graphed: \quad $y=$ \rule{4cm}{0.4pt} \qquad $y=$ \rule{4cm}{0.4pt}

\vspace{14pt}
The graphing calculator shows the curve and line intersect at $x\approx$ \rule{2cm}{0.4pt} and $x\approx$ \rule{2cm}{0.4pt}.

\vfill
\noindent\textbf{Scoring:} Questions 1 and 5 are 10 points each. Question 2 is a bonus of up to 5 points. Total out of 20.

\newpage
\noindent{\LARGE\textbf{2-7 Lesson Quiz, short form: answer key}}

\vspace{4pt}\noindent\rule{\textwidth}{0.6pt}

\begin{enumerate}
\item[\textbf{1.}] \textbf{A, 2.} $x^2-3=-2x+4$, so $x^2+2x-7=0$; the discriminant is $2^2-4(1)(-7)=32>0$, so two real solutions ($x=-1\pm2\sqrt2$, about $1.83$ and $-3.83$). A sketch shows the line crossing the parabola (vertex $(0,-3)$) twice. DOK 2, Realize skills A91, A54; IXL VQU (Example 1). \textbf{10 points:} 6 for A, 4 for a reason (a sketch with two crossings, or the discriminant).
\item[\textbf{2.}] \textbf{C, $(-4,-8)$ and $(2,-2)$.} $-\tfrac12x^2=x-4$, so $x^2+2x-8=0$, $(x+4)(x-2)=0$, $x=-4$ or $x=2$; then $y=x-4$ gives $-8$ and $-2$. Distractor A is the $x$-values without their $y$-values; B swaps the signs of $x$. DOK 2, Realize skills A91, A54; IXL HVZ (the Example 1 check; the packet's Exit question). \textbf{Bonus 5:} 5 for C with the substitution shown; 3 for C alone; 2 for A with correct work (the $x$-values found, not paired).
\item[\textbf{5.}] \textbf{$x\approx3.2$ and $x\approx-4.4$.} Graph $y=2x^2+5x-8$ and $y=\tfrac52x+20$; the intersections are about $(3.17,27.92)$ and $(-4.42,8.96)$. Check: $2x^2+\tfrac52x-28=0$ gives $x=\frac{-2.5\pm\sqrt{230.25}}{4}$, about $3.168$ and $-4.418$. DOK 2, Realize skills A91, A54; IXL VQU (Example 5). \textbf{10 points:} 2 for the system, 4 for each $x$-value to the nearest tenth (2 each if not rounded to tenths).
\end{enumerate}

\textbf{Items left off this form} (full Savvas quiz only; available for a retake or Form B): Savvas item 3, Nate tosses a ball up a hill, path $y=-\tfrac14x^2+\tfrac{33}{5}x$, hill $y=\tfrac15x$: how far does the ball travel horizontally (answer 25.6 ft; DOK 1, A91, A54; taught by Example 3, cut in the fast packet); Savvas item 4, solve the system of inequalities $y\geq2x^2+2x-3$, $y\leq\tfrac13x+6$ using shading (answer: both boundaries solid, shade above the parabola and below the line where they overlap; DOK 2, A91, A54; taught by Example 4, cut).

\textbf{Topic 2 test:} item 1 is the same task as Topic Assessment question 16, and item 5 the same task as question 17 (to the nearest hundredth there).

\textbf{IXL replacement:} SmartScore 70 or above on VQU and HVZ replaces this quiz with 20/20.

\end{document}
